% 完备空间（综述）
% license CCBYSA3
% type Wiki

本文根据 CC-BY-SA 协议转载翻译自维基百科\href{https://en.wikipedia.org/wiki/Complete_metric_space}{相关文章}

在数学分析中，如果一个度量空间$M$中的每个柯西序列都收敛于$M$中的某个点，那么这个度量空间$M$就称为完备的，也称为柯西空间。

直观地说，一个空间是完备的，意味着它内部或边界上不存在“缺失的点”。例如，有理数集并不是完备的，因为$\sqrt{2}$这样的数“缺失”在有理数集中，尽管我们可以构造一个收敛于$\sqrt{2}$的有理数柯西序列（见下文更多示例）。总是可以通过“填补所有这些空缺”来得到某个空间的完备化，具体方法将在下文解释。
\subsection{定义}
\textbf{柯西序列}

一个度量空间$(X, d)$中的序列$x_{1}, x_{2}, x_{3}, \ldots$称为柯西序列，如果对于任意一个正实数$r > 0$，存在一个正整数$N$，使得对于所有大于$N$的正整数 $m, n$，都有：
$$
d(x_{m}, x_{n}) < r.~
$$
\textbf{完备空间}

一个度量空间$(X, d)$称为完备的，当且仅当满足以下任一等价条件：
\begin{enumerate}
\item 柯西序列收敛性，$X$ 中的每一个柯西序列都在 $X$ 中收敛（即收敛到 $X$ 中的某个点）。
\item 嵌套闭集的交集不为空，$X$ 中任意一个非空闭集序列$\{F_n\}$，如果满足以下条件：$F_{n+1} \subseteq F_n$（即序列是逐步缩小的嵌套闭集序列）；集合的直径 $\operatorname{diam}(F_n)$ 趋于 0：$\operatorname{diam}(F_n) \to 0$，那么这些集合的交集$\bigcap_{n=1}^{\infty} F_n$恰好包含一个唯一的点 $x \in X$。
\end{enumerate}
\subsection{示例}
有理数集 $\mathbb{Q}$ 搭配标准度量（即差的绝对值）并不是完备的。例如，考虑如下定义的序列：
$$
x_1 = 1, \quad x_{n+1} = \frac{x_n}{2} + \frac{1}{x_n}.~
$$
这是一个由有理数组成的柯西序列，但它并不收敛于任何一个有理数。如果这个序列确实有极限 $x$，那么根据递推公式有：$x = \frac{x}{2} + \frac{1}{x}$，从而得到$
x^2 = 2$。然而，没有任何一个有理数满足这个条件。但如果把该序列放在实数集里来看，它确实收敛到无理数 $\sqrt{2}$。

开区间 $(0,1)$ 也不是完备的（度量依然是绝对值差）。考虑序列：$x_n = \frac{1}{n}$。这是一个柯西序列，但在 $(0,1)$ 中没有极限。然而，闭区间 $[0,1]$ 是完备的；同样的序列在这个空间中有极限，其极限点就是 $0$。
